\documentclass[10pt]{article}
\usepackage{tmlr}
%%%%% NEW MATH DEFINITIONS %%%%%

\usepackage{amsmath,amsfonts,bm}

% Mark sections of captions for referring to divisions of figures
\newcommand{\figleft}{{\em (Left)}}
\newcommand{\figcenter}{{\em (Center)}}
\newcommand{\figright}{{\em (Right)}}
\newcommand{\figtop}{{\em (Top)}}
\newcommand{\figbottom}{{\em (Bottom)}}
\newcommand{\captiona}{{\em (a)}}
\newcommand{\captionb}{{\em (b)}}
\newcommand{\captionc}{{\em (c)}}
\newcommand{\captiond}{{\em (d)}}

% Highlight a newly defined term
\newcommand{\newterm}[1]{{\bf #1}}


% Figure reference, lower-case.
\def\figref#1{figure~\ref{#1}}
% Figure reference, capital. For start of sentence
\def\Figref#1{Figure~\ref{#1}}
\def\twofigref#1#2{figures \ref{#1} and \ref{#2}}
\def\quadfigref#1#2#3#4{figures \ref{#1}, \ref{#2}, \ref{#3} and \ref{#4}}
% Section reference, lower-case.
\def\secref#1{section~\ref{#1}}
% Section reference, capital.
\def\Secref#1{Section~\ref{#1}}
% Reference to two sections.
\def\twosecrefs#1#2{sections \ref{#1} and \ref{#2}}
% Reference to three sections.
\def\secrefs#1#2#3{sections \ref{#1}, \ref{#2} and \ref{#3}}
% Reference to an equation, lower-case.
\def\eqref#1{equation~\ref{#1}}
% Reference to an equation, upper case
\def\Eqref#1{Equation~\ref{#1}}
% A raw reference to an equation---avoid using if possible
\def\plaineqref#1{\ref{#1}}
% Reference to a chapter, lower-case.
\def\chapref#1{chapter~\ref{#1}}
% Reference to an equation, upper case.
\def\Chapref#1{Chapter~\ref{#1}}
% Reference to a range of chapters
\def\rangechapref#1#2{chapters\ref{#1}--\ref{#2}}
% Reference to an algorithm, lower-case.
\def\algref#1{algorithm~\ref{#1}}
% Reference to an algorithm, upper case.
\def\Algref#1{Algorithm~\ref{#1}}
\def\twoalgref#1#2{algorithms \ref{#1} and \ref{#2}}
\def\Twoalgref#1#2{Algorithms \ref{#1} and \ref{#2}}
% Reference to a part, lower case
\def\partref#1{part~\ref{#1}}
% Reference to a part, upper case
\def\Partref#1{Part~\ref{#1}}
\def\twopartref#1#2{parts \ref{#1} and \ref{#2}}

\def\ceil#1{\lceil #1 \rceil}
\def\floor#1{\lfloor #1 \rfloor}
\def\1{\bm{1}}
\newcommand{\train}{\mathcal{D}}
\newcommand{\valid}{\mathcal{D_{\mathrm{valid}}}}
\newcommand{\test}{\mathcal{D_{\mathrm{test}}}}

\def\eps{{\epsilon}}


% Random variables
\def\reta{{\textnormal{$\eta$}}}
\def\ra{{\textnormal{a}}}
\def\rb{{\textnormal{b}}}
\def\rc{{\textnormal{c}}}
\def\rd{{\textnormal{d}}}
\def\re{{\textnormal{e}}}
\def\rf{{\textnormal{f}}}
\def\rg{{\textnormal{g}}}
\def\rh{{\textnormal{h}}}
\def\ri{{\textnormal{i}}}
\def\rj{{\textnormal{j}}}
\def\rk{{\textnormal{k}}}
\def\rl{{\textnormal{l}}}
% rm is already a command, just don't name any random variables m
\def\rn{{\textnormal{n}}}
\def\ro{{\textnormal{o}}}
\def\rp{{\textnormal{p}}}
\def\rq{{\textnormal{q}}}
\def\rr{{\textnormal{r}}}
\def\rs{{\textnormal{s}}}
\def\rt{{\textnormal{t}}}
\def\ru{{\textnormal{u}}}
\def\rv{{\textnormal{v}}}
\def\rw{{\textnormal{w}}}
\def\rx{{\textnormal{x}}}
\def\ry{{\textnormal{y}}}
\def\rz{{\textnormal{z}}}

% Random vectors
\def\rvepsilon{{\mathbf{\epsilon}}}
\def\rvtheta{{\mathbf{\theta}}}
\def\rva{{\mathbf{a}}}
\def\rvb{{\mathbf{b}}}
\def\rvc{{\mathbf{c}}}
\def\rvd{{\mathbf{d}}}
\def\rve{{\mathbf{e}}}
\def\rvf{{\mathbf{f}}}
\def\rvg{{\mathbf{g}}}
\def\rvh{{\mathbf{h}}}
\def\rvu{{\mathbf{i}}}
\def\rvj{{\mathbf{j}}}
\def\rvk{{\mathbf{k}}}
\def\rvl{{\mathbf{l}}}
\def\rvm{{\mathbf{m}}}
\def\rvn{{\mathbf{n}}}
\def\rvo{{\mathbf{o}}}
\def\rvp{{\mathbf{p}}}
\def\rvq{{\mathbf{q}}}
\def\rvr{{\mathbf{r}}}
\def\rvs{{\mathbf{s}}}
\def\rvt{{\mathbf{t}}}
\def\rvu{{\mathbf{u}}}
\def\rvv{{\mathbf{v}}}
\def\rvw{{\mathbf{w}}}
\def\rvx{{\mathbf{x}}}
\def\rvy{{\mathbf{y}}}
\def\rvz{{\mathbf{z}}}

% Elements of random vectors
\def\erva{{\textnormal{a}}}
\def\ervb{{\textnormal{b}}}
\def\ervc{{\textnormal{c}}}
\def\ervd{{\textnormal{d}}}
\def\erve{{\textnormal{e}}}
\def\ervf{{\textnormal{f}}}
\def\ervg{{\textnormal{g}}}
\def\ervh{{\textnormal{h}}}
\def\ervi{{\textnormal{i}}}
\def\ervj{{\textnormal{j}}}
\def\ervk{{\textnormal{k}}}
\def\ervl{{\textnormal{l}}}
\def\ervm{{\textnormal{m}}}
\def\ervn{{\textnormal{n}}}
\def\ervo{{\textnormal{o}}}
\def\ervp{{\textnormal{p}}}
\def\ervq{{\textnormal{q}}}
\def\ervr{{\textnormal{r}}}
\def\ervs{{\textnormal{s}}}
\def\ervt{{\textnormal{t}}}
\def\ervu{{\textnormal{u}}}
\def\ervv{{\textnormal{v}}}
\def\ervw{{\textnormal{w}}}
\def\ervx{{\textnormal{x}}}
\def\ervy{{\textnormal{y}}}
\def\ervz{{\textnormal{z}}}

% Random matrices
\def\rmA{{\mathbf{A}}}
\def\rmB{{\mathbf{B}}}
\def\rmC{{\mathbf{C}}}
\def\rmD{{\mathbf{D}}}
\def\rmE{{\mathbf{E}}}
\def\rmF{{\mathbf{F}}}
\def\rmG{{\mathbf{G}}}
\def\rmH{{\mathbf{H}}}
\def\rmI{{\mathbf{I}}}
\def\rmJ{{\mathbf{J}}}
\def\rmK{{\mathbf{K}}}
\def\rmL{{\mathbf{L}}}
\def\rmM{{\mathbf{M}}}
\def\rmN{{\mathbf{N}}}
\def\rmO{{\mathbf{O}}}
\def\rmP{{\mathbf{P}}}
\def\rmQ{{\mathbf{Q}}}
\def\rmR{{\mathbf{R}}}
\def\rmS{{\mathbf{S}}}
\def\rmT{{\mathbf{T}}}
\def\rmU{{\mathbf{U}}}
\def\rmV{{\mathbf{V}}}
\def\rmW{{\mathbf{W}}}
\def\rmX{{\mathbf{X}}}
\def\rmY{{\mathbf{Y}}}
\def\rmZ{{\mathbf{Z}}}

% Elements of random matrices
\def\ermA{{\textnormal{A}}}
\def\ermB{{\textnormal{B}}}
\def\ermC{{\textnormal{C}}}
\def\ermD{{\textnormal{D}}}
\def\ermE{{\textnormal{E}}}
\def\ermF{{\textnormal{F}}}
\def\ermG{{\textnormal{G}}}
\def\ermH{{\textnormal{H}}}
\def\ermI{{\textnormal{I}}}
\def\ermJ{{\textnormal{J}}}
\def\ermK{{\textnormal{K}}}
\def\ermL{{\textnormal{L}}}
\def\ermM{{\textnormal{M}}}
\def\ermN{{\textnormal{N}}}
\def\ermO{{\textnormal{O}}}
\def\ermP{{\textnormal{P}}}
\def\ermQ{{\textnormal{Q}}}
\def\ermR{{\textnormal{R}}}
\def\ermS{{\textnormal{S}}}
\def\ermT{{\textnormal{T}}}
\def\ermU{{\textnormal{U}}}
\def\ermV{{\textnormal{V}}}
\def\ermW{{\textnormal{W}}}
\def\ermX{{\textnormal{X}}}
\def\ermY{{\textnormal{Y}}}
\def\ermZ{{\textnormal{Z}}}

% Vectors
\def\vzero{{\bm{0}}}
\def\vone{{\bm{1}}}
\def\vmu{{\bm{\mu}}}
\def\vtheta{{\bm{\theta}}}
\def\va{{\bm{a}}}
\def\vb{{\bm{b}}}
\def\vc{{\bm{c}}}
\def\vd{{\bm{d}}}
\def\ve{{\bm{e}}}
\def\vf{{\bm{f}}}
\def\vg{{\bm{g}}}
\def\vh{{\bm{h}}}
\def\vi{{\bm{i}}}
\def\vj{{\bm{j}}}
\def\vk{{\bm{k}}}
\def\vl{{\bm{l}}}
\def\vm{{\bm{m}}}
\def\vn{{\bm{n}}}
\def\vo{{\bm{o}}}
\def\vp{{\bm{p}}}
\def\vq{{\bm{q}}}
\def\vr{{\bm{r}}}
\def\vs{{\bm{s}}}
\def\vt{{\bm{t}}}
\def\vu{{\bm{u}}}
\def\vv{{\bm{v}}}
\def\vw{{\bm{w}}}
\def\vx{{\bm{x}}}
\def\vy{{\bm{y}}}
\def\vz{{\bm{z}}}

% Elements of vectors
\def\evalpha{{\alpha}}
\def\evbeta{{\beta}}
\def\evepsilon{{\epsilon}}
\def\evlambda{{\lambda}}
\def\evomega{{\omega}}
\def\evmu{{\mu}}
\def\evpsi{{\psi}}
\def\evsigma{{\sigma}}
\def\evtheta{{\theta}}
\def\eva{{a}}
\def\evb{{b}}
\def\evc{{c}}
\def\evd{{d}}
\def\eve{{e}}
\def\evf{{f}}
\def\evg{{g}}
\def\evh{{h}}
\def\evi{{i}}
\def\evj{{j}}
\def\evk{{k}}
\def\evl{{l}}
\def\evm{{m}}
\def\evn{{n}}
\def\evo{{o}}
\def\evp{{p}}
\def\evq{{q}}
\def\evr{{r}}
\def\evs{{s}}
\def\evt{{t}}
\def\evu{{u}}
\def\evv{{v}}
\def\evw{{w}}
\def\evx{{x}}
\def\evy{{y}}
\def\evz{{z}}

% Matrix
\def\mA{{\bm{A}}}
\def\mB{{\bm{B}}}
\def\mC{{\bm{C}}}
\def\mD{{\bm{D}}}
\def\mE{{\bm{E}}}
\def\mF{{\bm{F}}}
\def\mG{{\bm{G}}}
\def\mH{{\bm{H}}}
\def\mI{{\bm{I}}}
\def\mJ{{\bm{J}}}
\def\mK{{\bm{K}}}
\def\mL{{\bm{L}}}
\def\mM{{\bm{M}}}
\def\mN{{\bm{N}}}
\def\mO{{\bm{O}}}
\def\mP{{\bm{P}}}
\def\mQ{{\bm{Q}}}
\def\mR{{\bm{R}}}
\def\mS{{\bm{S}}}
\def\mT{{\bm{T}}}
\def\mU{{\bm{U}}}
\def\mV{{\bm{V}}}
\def\mW{{\bm{W}}}
\def\mX{{\bm{X}}}
\def\mY{{\bm{Y}}}
\def\mZ{{\bm{Z}}}
\def\mBeta{{\bm{\beta}}}
\def\mPhi{{\bm{\Phi}}}
\def\mLambda{{\bm{\Lambda}}}
\def\mSigma{{\bm{\Sigma}}}

% Tensor
\DeclareMathAlphabet{\mathsfit}{\encodingdefault}{\sfdefault}{m}{sl}
\SetMathAlphabet{\mathsfit}{bold}{\encodingdefault}{\sfdefault}{bx}{n}
\newcommand{\tens}[1]{\bm{\mathsfit{#1}}}
\def\tA{{\tens{A}}}
\def\tB{{\tens{B}}}
\def\tC{{\tens{C}}}
\def\tD{{\tens{D}}}
\def\tE{{\tens{E}}}
\def\tF{{\tens{F}}}
\def\tG{{\tens{G}}}
\def\tH{{\tens{H}}}
\def\tI{{\tens{I}}}
\def\tJ{{\tens{J}}}
\def\tK{{\tens{K}}}
\def\tL{{\tens{L}}}
\def\tM{{\tens{M}}}
\def\tN{{\tens{N}}}
\def\tO{{\tens{O}}}
\def\tP{{\tens{P}}}
\def\tQ{{\tens{Q}}}
\def\tR{{\tens{R}}}
\def\tS{{\tens{S}}}
\def\tT{{\tens{T}}}
\def\tU{{\tens{U}}}
\def\tV{{\tens{V}}}
\def\tW{{\tens{W}}}
\def\tX{{\tens{X}}}
\def\tY{{\tens{Y}}}
\def\tZ{{\tens{Z}}}


% Graph
\def\gA{{\mathcal{A}}}
\def\gB{{\mathcal{B}}}
\def\gC{{\mathcal{C}}}
\def\gD{{\mathcal{D}}}
\def\gE{{\mathcal{E}}}
\def\gF{{\mathcal{F}}}
\def\gG{{\mathcal{G}}}
\def\gH{{\mathcal{H}}}
\def\gI{{\mathcal{I}}}
\def\gJ{{\mathcal{J}}}
\def\gK{{\mathcal{K}}}
\def\gL{{\mathcal{L}}}
\def\gM{{\mathcal{M}}}
\def\gN{{\mathcal{N}}}
\def\gO{{\mathcal{O}}}
\def\gP{{\mathcal{P}}}
\def\gQ{{\mathcal{Q}}}
\def\gR{{\mathcal{R}}}
\def\gS{{\mathcal{S}}}
\def\gT{{\mathcal{T}}}
\def\gU{{\mathcal{U}}}
\def\gV{{\mathcal{V}}}
\def\gW{{\mathcal{W}}}
\def\gX{{\mathcal{X}}}
\def\gY{{\mathcal{Y}}}
\def\gZ{{\mathcal{Z}}}

% Sets
\def\sA{{\mathbb{A}}}
\def\sB{{\mathbb{B}}}
\def\sC{{\mathbb{C}}}
\def\sD{{\mathbb{D}}}
% Don't use a set called E, because this would be the same as our symbol
% for expectation.
\def\sF{{\mathbb{F}}}
\def\sG{{\mathbb{G}}}
\def\sH{{\mathbb{H}}}
\def\sI{{\mathbb{I}}}
\def\sJ{{\mathbb{J}}}
\def\sK{{\mathbb{K}}}
\def\sL{{\mathbb{L}}}
\def\sM{{\mathbb{M}}}
\def\sN{{\mathbb{N}}}
\def\sO{{\mathbb{O}}}
\def\sP{{\mathbb{P}}}
\def\sQ{{\mathbb{Q}}}
\def\sR{{\mathbb{R}}}
\def\sS{{\mathbb{S}}}
\def\sT{{\mathbb{T}}}
\def\sU{{\mathbb{U}}}
\def\sV{{\mathbb{V}}}
\def\sW{{\mathbb{W}}}
\def\sX{{\mathbb{X}}}
\def\sY{{\mathbb{Y}}}
\def\sZ{{\mathbb{Z}}}

% Entries of a matrix
\def\emLambda{{\Lambda}}
\def\emA{{A}}
\def\emB{{B}}
\def\emC{{C}}
\def\emD{{D}}
\def\emE{{E}}
\def\emF{{F}}
\def\emG{{G}}
\def\emH{{H}}
\def\emI{{I}}
\def\emJ{{J}}
\def\emK{{K}}
\def\emL{{L}}
\def\emM{{M}}
\def\emN{{N}}
\def\emO{{O}}
\def\emP{{P}}
\def\emQ{{Q}}
\def\emR{{R}}
\def\emS{{S}}
\def\emT{{T}}
\def\emU{{U}}
\def\emV{{V}}
\def\emW{{W}}
\def\emX{{X}}
\def\emY{{Y}}
\def\emZ{{Z}}
\def\emSigma{{\Sigma}}

% entries of a tensor
% Same font as tensor, without \bm wrapper
\newcommand{\etens}[1]{\mathsfit{#1}}
\def\etLambda{{\etens{\Lambda}}}
\def\etA{{\etens{A}}}
\def\etB{{\etens{B}}}
\def\etC{{\etens{C}}}
\def\etD{{\etens{D}}}
\def\etE{{\etens{E}}}
\def\etF{{\etens{F}}}
\def\etG{{\etens{G}}}
\def\etH{{\etens{H}}}
\def\etI{{\etens{I}}}
\def\etJ{{\etens{J}}}
\def\etK{{\etens{K}}}
\def\etL{{\etens{L}}}
\def\etM{{\etens{M}}}
\def\etN{{\etens{N}}}
\def\etO{{\etens{O}}}
\def\etP{{\etens{P}}}
\def\etQ{{\etens{Q}}}
\def\etR{{\etens{R}}}
\def\etS{{\etens{S}}}
\def\etT{{\etens{T}}}
\def\etU{{\etens{U}}}
\def\etV{{\etens{V}}}
\def\etW{{\etens{W}}}
\def\etX{{\etens{X}}}
\def\etY{{\etens{Y}}}
\def\etZ{{\etens{Z}}}

% The true underlying data generating distribution
\newcommand{\pdata}{p_{\rm{data}}}
% The empirical distribution defined by the training set
\newcommand{\ptrain}{\hat{p}_{\rm{data}}}
\newcommand{\Ptrain}{\hat{P}_{\rm{data}}}
% The model distribution
\newcommand{\pmodel}{p_{\rm{model}}}
\newcommand{\Pmodel}{P_{\rm{model}}}
\newcommand{\ptildemodel}{\tilde{p}_{\rm{model}}}
% Stochastic autoencoder distributions
\newcommand{\pencode}{p_{\rm{encoder}}}
\newcommand{\pdecode}{p_{\rm{decoder}}}
\newcommand{\precons}{p_{\rm{reconstruct}}}

\newcommand{\laplace}{\mathrm{Laplace}} % Laplace distribution

\newcommand{\E}{\mathbb{E}}
\newcommand{\Ls}{\mathcal{L}}
\newcommand{\R}{\mathbb{R}}
\newcommand{\emp}{\tilde{p}}
\newcommand{\lr}{\alpha}
\newcommand{\reg}{\lambda}
\newcommand{\rect}{\mathrm{rectifier}}
\newcommand{\softmax}{\mathrm{softmax}}
\newcommand{\sigmoid}{\sigma}
\newcommand{\softplus}{\zeta}
\newcommand{\KL}{D_{\mathrm{KL}}}
\newcommand{\Var}{\mathrm{Var}}
\newcommand{\standarderror}{\mathrm{SE}}
\newcommand{\Cov}{\mathrm{Cov}}
% Wolfram Mathworld says $L^2$ is for function spaces and $\ell^2$ is for vectors
% But then they seem to use $L^2$ for vectors throughout the site, and so does
% wikipedia.
\newcommand{\normlzero}{L^0}
\newcommand{\normlone}{L^1}
\newcommand{\normltwo}{L^2}
\newcommand{\normlp}{L^p}
\newcommand{\normmax}{L^\infty}

\newcommand{\parents}{Pa} % See usage in notation.tex. Chosen to match Daphne's book.

\DeclareMathOperator*{\argmax}{arg\,max}
\DeclareMathOperator*{\argmin}{arg\,min}

\DeclareMathOperator{\sign}{sign}
\DeclareMathOperator{\Tr}{Tr}
\let\ab\allowbreak

\usepackage{hyperref}
\usepackage{url}
\usepackage{longtable}
\usepackage{booktabs}
\usepackage{array}
\usepackage{calc}
\usepackage{newunicodechar}
\providecommand{\tightlist}{\setlength{\itemsep}{0pt}\setlength{\parskip}{0pt}}
\newunicodechar{σ}{\ensuremath{\sigma}}
\newunicodechar{±}{\ensuremath{\pm}}
\newunicodechar{→}{\ensuremath{\rightarrow}}
\newunicodechar{−}{\ensuremath{-}}
\newunicodechar{×}{\ensuremath{\times}}
\newunicodechar{·}{\ensuremath{\cdot}}
\newunicodechar{α}{\ensuremath{\alpha}}
\newunicodechar{ρ}{\ensuremath{\rho}}
\newunicodechar{Σ}{\ensuremath{\Sigma}}
\newunicodechar{≥}{\ensuremath{\geq}}
\newunicodechar{≈}{\ensuremath{\approx}}

\title{Measured, Not Asked: Free Salience Signals for Consolidation and Routing at Small Scale}
\author{Anonymous authors}

\begin{document}
\maketitle

\begin{abstract}
An agent that moves knowledge from its context into its weights, or decides which part of a
long context to read, needs a salience signal. Three such signals come free, with no reward
model or judge: the keep/drop decision a context compactor already makes, the divergence
between a model's predictions with and without a document in context, and the region a model
declares it needs to attend to. We test all three on Qwen2.5 models from 0.5B to 7B with a
matched control beside every signal-strength figure, and report two findings. \textbf{First,
measured beats asked.} Signals the model is asked to state fail or are weak: asked to list
spans worth keeping, 0.5B and 1.5B models return the prefix \texttt{0,\ 1,\ 2,\ ...} on every event;
asked to declare which of eight regions holds an answer, the 1.5B model does no better than
naming a fixed region. The same models' measured signals work: reading the keep decision off
the logits gives a 2.56x salience lift, and scoring each region by the log-odds of a
relevance question routes at 0.33--0.42 hit rate against 0.125 chance, beating the stated
declaration at 0.5B, 1.5B and 7B (by 0.19 at 7B, t(2)=24) and late-layer attention probing
at 0.5B and 7B. \textbf{Second, a salient signal is not necessarily a good training target.}
Consolidating on the compactor's keep decision, a +15.85σ salience signal, recovers fewer
evicted facts than consolidating on uniformly sampled spans on all five synthetic corpora
(1.7\% vs 7.6\%, paired t(4)=3.10), because the compactor has systematic blind spots; and
gating distillation on the context gap performs no better than random token selection under
identical masking (5 seeds, t(4)=0.74). Neither free signal beats uncurated coverage as a
training target. Matched controls caught eight artifacts in our own pipeline, six of which
had inflated a result in our favour.
\end{abstract}

\section{Introduction}

Context is expensive and transient; weights are cheap to read and permanent. A long-running
agent therefore wants to move what matters from one to the other, and to read only the part
of its context a question needs. In both cases the training or routing step is routine --- a
LoRA pass, a sparse KV read --- and the hard part is deciding \emph{what} matters.

Prior work pays for that decision. Nightly consolidation \citep{dennis2026beyond} writes reflections with a second
generation pass; SCoL \citep{wang2026scol} learns where to write with meta-RL; CompactionRL \citep{li2026compactionrl} trains the
compactor with RL. We ask what is already free, and find three candidates:

\begin{itemize}
\tightlist
\item
  \textbf{The compaction decision.} Every time an agent compacts its context, it decides which
  spans survive verbatim. That is a salience label, emitted at no extra cost.
\item
  \textbf{The context gap.} Run a frozen model on the same continuation with and without a skill
  document in context. The per-token KL divergence between the two is the document's causal
  contribution, from two forward passes.
\item
  \textbf{The attention declaration.} Declarative Attention \citep{ho2026attention} has a model state which region of
  its context it needs, and skips the rest of the KV-cache read. The declaration is also a
  salience label.
\end{itemize}

We test each as a training or routing signal on models a single GPU can run, 0.5B to 7B. The
results sort along two axes. The first is how the signal is elicited: \emph{measured} from the
model's logits or activations, or \emph{asked} of it as generated text. This decides whether a
signal exists at small scale. The second is whether a signal that exists is the right thing
to train on. This decides whether it helps.

\textbf{Contributions.} This is a measurement paper: the methods are hypotheses under test, and
two of the three consolidation hypotheses come back negative.

\begin{enumerate}
\def\labelenumi{\arabic{enumi}.}
\tightlist
\item
  \textbf{Measured beats asked} (§4). On three independent mechanisms, a signal the model is
  asked to state fails or is weak, while the same model's measured signal works. A generated
  span list is exactly random; the same keep decision read off the logits gives a 2.56x lift.
  A generated attention declaration is at a fixed-slot control at 1.5B and carries under half
  the content signal of a logit read at 0.5B and 7B.
\item
  \textbf{Salient is not trainable} (§5). Two free signals that clear matched controls do not
  beat uniform coverage as training targets. Compaction-gated consolidation is significantly
  worse than uniform replay across five corpora; context-gap gating is indistinguishable from
  random token selection. They fail for different reasons --- coverage and neutrality.
\item
  \textbf{A logit-read substitute for declarative attention below the scale where declaring
  works} (§4.2): a cheap relevance query that routes at about 3x chance at every scale
  tested, with a content-shuffle control that separates content-driven from position-driven
  routing.
\item
  \textbf{A matched control beside every signal-strength figure} (§6), which caught eight
  artifacts in our own pipeline.
\end{enumerate}

\section{Related work}

\textbf{Consolidation from context.} Beyond Inference-Only Deployment \citep{dennis2026beyond} shows that periodic LoRA
consolidation beats cascading compaction, and establishes median rather than mean per-token
cross-entropy as the metric that tracks judged accuracy (r = +0.99 against −0.51). Its
targets are reflected facts written by a second generation pass; we consolidate on the
compactor's raw decision instead, which removes that pass and reflection quality as a
confound. Their recipe is our baseline. SCoL \citep{wang2026scol} learns \emph{where} to write consolidated
knowledge --- the complementary axis to our \emph{what}. What to Keep, What to Forget \citep{colaco2026keep} frames
compaction as rate-distortion and stops at the compactor; we consume its decisions as
training signal. CompactionRL \citep{li2026compactionrl} trains the compactor with RL; we hold it frozen and train
the backbone, and the two compose. Compression-Aware Abstention \citep{khodabandehlou2026abstention} trains a model on
KV-compression masks to refuse when compression removed the answer evidence; we reproduce its direction at
0.5B (§5.1).

\textbf{Importance-gated distillation and token selection.} Skill-to-LoRA \citep{zhang2026skill} distils each skill
into its own adapter by uniform self-distillation, and is our baseline in §5.2. ThinkSwitch
\citep{saini2026thinkswitch} distils context into LoRA adapters with weight interpolation, also ungated. PEAM \citep{guo2026peam}
internalises an embodied agent's episodic experience into its parameters by contrasting
trajectories --- the nearest neighbour to our gate, on different objects with a different
signal. Titans \citep{behrouz2024titans} uses the gradient of an
associative-memory loss as surprise at test time; our context gap is a surprise signal in the
same spirit, computed as a divergence between two forward passes and applied offline. Rho-1\citep{lin2024rho} and TIP \citep{xu2026tip} report the opposite of our §5.2 result in their settings: selective loss on
scored tokens helps pretraining \citep{lin2024rho}, and training on under 10\% of tokens, chosen by
teacher--student divergence and entropy, nearly matches full-token on-policy distillation
\citep{xu2026tip}. Our design isolates the ranking from the masking with a random-subset control at
identical masking. In our setting the subsetting carries the benefit and the ranking adds
nothing; whether that decomposition holds in theirs is open.

\textbf{Attention routing.} Declarative Attention \citep{ho2026attention} has a 27--31B model declare \texttt{\textless{}global\textgreater{}},
\texttt{\textless{}focus\textgreater{}} or \texttt{\textless{}local\textgreater{}} in its chain of thought and skips most of the KV-cache read
accordingly: 52.0\% and 31.1\% fewer attended tokens at −1.27 and −2.75 points of accuracy,
zero-shot, with accuracy losses that ``shrink with model scale''. We test the declaration from
below, at the scales a single GPU runs. Sentinel \citep{zhang2025sentinel} and the retrieval-head line \citep{wu2024retrieval} route
by probing attention instead of asking; we run that route on the same models, layouts and
probes as the other two.

\section{Setup and measurement protocol}

\textbf{Models and compute.} Qwen2.5-Instruct \citep{qwen2024qwen25} at 0.5B, 1.5B and 7B; LoRA \citep{hu2021lora} rank 16,
α = 32, on all attention and MLP projections. 0.5B and 1.5B run in fp16; 7B runs in nf4 4-bit
quantisation to fit a single 16 GB T4, which is all the compute used.

\textbf{Salience lift.} For a signal that marks spans, lift is the keep rate of fact spans over
the keep rate of filler spans. Every lift is reported against a \textbf{positional control} ---
what keeping the first N spans scores on the same data --- and with its σ over chance and over
that control.

\textbf{Statistics.} Single-run comparisons use a binomial σ. Multi-seed comparisons use a paired
t across seeds against the two-sided 5\% critical value for k−1 degrees of freedom (4.30 at
three seeds, 2.78 at five). We call a margin significant only if it clears that value; a
difference that does not is reported as a tie, not a trend.

\textbf{Content shuffle} (§4.2). For routing signals, a raw hit rate cannot distinguish tracking
content from favouring a position. We permute which content sits in which region and
re-elicit. \texttt{content\_dependence} is the shuffled hit rate above chance; \texttt{slot\_stable\_rate} is
how often the choice stays put when the content moves (near 1/K is content-driven, near 1
position-driven). A best-constant control --- always naming the most common gold region ---
exposes a model that names one region regardless.

\section{Finding 1: measured beats asked}

\subsection{The compaction decision}

The compactor segments the evictable prefix of a trajectory into sentence spans and keeps
\texttt{ceil(ρ·n)} of them, ρ = 0.25. We elicit the keep decision two ways: \textbf{generate}, where the
model lists span indices to keep; and \textbf{read}, where the model is asked of each span whether
to keep it and we score \texttt{log\ p(Yes)\ −\ log\ p(No)} off the logits, keeping the top \texttt{ρ·n}. Spans
are shuffled before presentation so position cannot masquerade as salience.

\begin{longtable}[]{@{}llll@{}}
\toprule\noalign{}
Elicitation & Backbone & Lift & Verdict \\
\midrule\noalign{}
\endhead
\bottomrule\noalign{}
\endlastfoot
generated index list & Qwen2.5-0.5B / 1.5B & 0.53x & below positional control \\
logit read & SmolLM2-360M & 2.17x & clears \\
logit read & Qwen2.5-0.5B & 0.42x & fails \\
logit read & Qwen2.5-1.5B & 2.56x & clears \\
\end{longtable}

Asked to list spans, both Qwen models return the contiguous prefix \texttt{0,\ 1,\ 2,\ ...} on 100\% of
events regardless of content --- under shuffling, exactly random selection. Read off the
logits, the 1.5B model's decision is a strong signal. On the corrected synthetic benchmark
used for §5.1 (48 trajectories, 97 compaction events, 6,244 spans), it keeps fact spans at
0.663 and filler at 0.236: a 2.81x lift, \textbf{+15.85σ over chance and +8.40σ over the positional
control}, with zero fallbacks or empty keeps. Per fact type it is uneven: \texttt{db\_engine},
\texttt{owner} and \texttt{version\_pin} are kept on 100\% of events, \texttt{config\_flag} on 3.4\% and
\texttt{auth\_header} on 0\%. §5.1 shows why that matters.

Measuring is necessary but not sufficient for a signal to exist: the 0.5B logit read fails
(0.42x) where SmolLM2-360M clears (2.17x) on identical data and code. Whether a model carries
a usable signal is a property of its behaviour on the probe, not of its size.

\subsection{The attention declaration}

\textbf{Task.} A HotpotQA trajectory with its supporting paragraphs scattered across the context
is split into K = 8 labelled regions, and the model gets a question whose answer lies in
exactly one. Rendered contexts reach 8,966 tokens, inside every model's 32,768-token window.
384 probes; the gold region is permuted per probe; three seeds reseed the layouts and the
shuffle over a fixed probe set, and the gold distribution for each seed is pinned so every
scale sees identical layouts.

\textbf{Elicitations.}

\begin{itemize}
\tightlist
\item
  \textbf{generate} --- the model replies \texttt{FOCUS:\ \textless{}k\textgreater{}}, parsed like a tool call, as in \citet{ho2026attention}.
\item
  \textbf{read} --- for each region we ask whether it contains the answer and score
  \texttt{log\ p(Yes)\ −\ log\ p(No)}; the argmax is the choice. K batched region prompts.
\item
  \textbf{attention} --- no elicitation: we run the model on the full context and question and take
  the attention mass the final query position places on each region, summed over layers.
  \textbf{attention\_late} restricts the sum to the second half of the layers. One forward pass
  over roughly the same tokens as \texttt{read}.
\end{itemize}

\textbf{Results} (mean ± sd over three seeds; chance 0.125, best-constant control 0.148):

\begin{longtable}[]{@{}lllllll@{}}
\toprule\noalign{}
model & elicitation & hit & σ over constant & content dependence & slot stable & unparsed \\
\midrule\noalign{}
\endhead
\bottomrule\noalign{}
\endlastfoot
0.5B & generate & 0.193 ± 0.012 & 2.2 & 0.085 ± 0.021 & 0.281 & 0.019 \\
0.5B & attention & 0.151 ± 0.003 & 0.2 & 0.014 ± 0.012 & 0.882 & 0.000 \\
0.5B & attention\_late & 0.234 ± 0.011 & 4.0 & 0.113 ± 0.029 & 0.356 & 0.000 \\
0.5B & \textbf{read} & \textbf{0.332 ± 0.004} & \textbf{7.7} & \textbf{0.203 ± 0.003} & 0.136 & 0.000 \\
1.5B & generate & 0.150 ± 0.008 & 0.1 & 0.030 ± 0.023 & 0.413 & 0.170 \\
1.5B & attention & 0.286 ± 0.023 & 6.0 & 0.163 ± 0.010 & 0.447 & 0.000 \\
1.5B & attention\_late & 0.374 ± 0.017 & 9.2 & 0.258 ± 0.028 & 0.200 & 0.000 \\
1.5B & \textbf{read} & \textbf{0.378 ± 0.005} & \textbf{9.3} & \textbf{0.241 ± 0.005} & 0.136 & 0.000 \\
7B & generate & 0.230 ± 0.010 & 3.8 & 0.116 ± 0.020 & 0.260 & 0.000 \\
7B & attention & 0.229 ± 0.009 & 3.8 & 0.085 ± 0.018 & 0.595 & 0.000 \\
7B & attention\_late & 0.359 ± 0.017 & 8.6 & 0.223 ± 0.019 & 0.242 & 0.000 \\
7B & \textbf{read} & \textbf{0.418 ± 0.011} & \textbf{10.7} & \textbf{0.285 ± 0.004} & 0.135 & 0.000 \\
\end{longtable}

\textbf{The stated declaration never approaches a measured one.} At 0.5B its raw hit rate clears
the constant control (+2.2σ), but the shuffle shows most of that is position: content
dependence 0.085, under half of \texttt{read}'s. At 1.5B it is indistinguishable from naming one
fixed region (+0.1σ), refuses the \texttt{FOCUS:} format on 17\% of probes, and its slot-stable rate
rises to 0.41. At 7B the refusals disappear and it clears the constant control again (+3.8σ,
between +3.7σ and +4.0σ on every seed), with content dependence 0.116 --- the 0.5B pattern
rather than the 1.5B collapse. \texttt{read} still beats it by 0.188 in hit rate (t(2) = 24.0) and
0.168 in content dependence (t(2) = 12.3). The decline from 0.5B to 1.5B does not continue;
the gap to the measured routes does.

\textbf{\texttt{read} wins at 0.5B and 7B and ties at 1.5B.} Against the best attention variant, \texttt{read}
leads at 0.5B on hit rate (0.332 vs 0.234, t(2) = 24.9) and content dependence (0.203 vs
0.113, t(2) = 4.91). At 1.5B the two tie: 0.378 vs 0.374 (t(2) = 0.30) and 0.241 vs 0.258
(t(2) = −0.85). At 7B \texttt{read} leads again, 0.418 vs 0.359 (t(2) = 12.0) and 0.285 vs 0.223
(t(2) = 4.69). \texttt{read} is the only signal whose content dependence rises at every scale
(0.203, 0.241, 0.285); \texttt{attention\_late} rises to 1.5B and falls back at 7B (0.113, 0.258,
0.223), so attention probing does not keep closing the gap. A single-seed run had shown
attention nominally ahead at 1.5B; three seeds showed that was layout noise.

\textbf{Layer choice decides whether probing works at all.} Summed over every layer, attention is
useless at 0.5B: content dependence 0.014, the same slot on 88\% of probes. Restricted to the
late half it clears its control at every scale. Early layers carry position and dominate the
sum --- one Qwen2.5-1.5B layer-0 query·key product reaches 152,967 against a late-layer peak
near 300. A probing baseline reported without this ablation would understate itself by more
than an order of magnitude.

\subsection{What the two mechanisms share}

On both the compaction decision and the attention declaration, the model's measured judgment
is better than its stated one --- the generated \texttt{FOCUS:} at 1.5B carries no content signal
(0.030) while the same model's logit read carries 0.241 and its late-layer attention 0.258.
The failure is not that small models lack the judgment; it is that they express it far less
well as structured output on demand than their logits reveal it, and scale to 7B narrows that
only partly. Measuring is also a design space rather than a single method: the self-query and
attention probing swap between win and tie across scales, and attention works only from late
layers. Of the routes we tested, the self-query is the safer default at every scale. The
context gap (§5.2) is measured by construction and needs no elicitation at all.

\section{Finding 2: salient is not trainable}

\subsection{Consolidating on the compaction decision}

\textbf{Method.} After every fourth compaction event, a sleep phase runs LoRA SFT on the spans the
compactor kept, with a reservoir replay buffer across phases. The compactor is Qwen2.5-1.5B
with the logit-read backend; the consolidated and evaluated backbone is Qwen2.5-0.5B.

\textbf{Benchmark.} Synthetic agent trajectories of 120 turns, each carrying six planted facts
(\texttt{auth\_header}, \texttt{db\_engine}, \ldots) in its first twelve turns and filler after, so compaction
evicts some facts. Each fact gets one probe, 288 per corpus over 48 trajectories. A probe is
correct if the normalised gold answer or an alias appears in the output, and \emph{evicted} if the
answer is absent from the compacted context --- the only probes where consolidation can help.
Every question names a project unique to its trajectory (§6 explains why this matters).

\textbf{Arms.} (a) uniform replay: the same sleep phases and budget on uniformly sampled spans,
kept and dropped alike; (b) reflection, following \citet{dennis2026beyond}; (c) cascading compaction with no
weight update; (d) full context, the ceiling; and ours, compaction-supervised.

\begin{longtable}[]{@{}lll@{}}
\toprule\noalign{}
method (corpus 0, 114 evicted) & all-probe & evicted \\
\midrule\noalign{}
\endhead
\bottomrule\noalign{}
\endlastfoot
no context (floor) & 0.000 & 0.000 \\
(c) cascading, no adapter & 0.413 & 0.000 \\
ours: compaction-supervised & 0.451 & \textbf{0.000} \\
ours + mask head & 0.323 & 0.009 \\
(b) reflection & 0.389 & 0.000 \\
\textbf{(a) uniform replay} & 0.465 & \textbf{0.123} \\
(d) full context & 0.712 & --- \\
\end{longtable}

Uniform replay recovers 14 of 114 evicted facts (+4.0σ over zero); compaction-gated
consolidation recovers none. Across five independently generated corpora:

\begin{longtable}[]{@{}lllll@{}}
\toprule\noalign{}
corpus & evicted & cascading & compaction & uniform \\
\midrule\noalign{}
\endhead
\bottomrule\noalign{}
\endlastfoot
0 & 114 & 0.000 & 0.000 & \textbf{0.123} \\
1 & 102 & 0.000 & 0.029 & \textbf{0.078} \\
2 & 112 & 0.000 & 0.000 & \textbf{0.009} \\
3 & 98 & 0.000 & 0.020 & \textbf{0.061} \\
4 & 112 & 0.000 & 0.036 & \textbf{0.107} \\
mean & & 0.000 & 0.017 & \textbf{0.076} \\
\end{longtable}

\textbf{Uniform beats the compaction gate on every corpus}, paired t(4) = 3.10 against a critical
value of 2.78. Corpus 0 was the most favourable draw; at three corpora the margin was not yet
significant (t(2) = 1.81). The gain is specific to evicted facts: all-probe accuracy is level
(t(4) = −0.88).

\textbf{Mechanism: coverage.} Both adapters fit their targets equally well (validation CE 0.0003
for ours, 0.0005 for uniform); they differ in what they fit. The compactor never keeps
\texttt{auth\_header} and almost never \texttt{config\_flag} (§4.1), so the gated adapter never sees them.
Uniform sampling covers them. A strong salience signal with a blind spot passes the blind spot
to the weights.

\textbf{Remixing dropped spans does not fix it.} Replacing 25\% or 50\% of the gated budget with
randomly drawn dropped spans, at a matched total budget, recovers 2 and 1 of 114. The dropped
pool is dominated by filler (5,956 filler spans), so a random draw rarely lands on the
excluded facts. Drawing round-robin across fact keys instead roughly doubles recovery at 25\%
(4 of 114) and falls to zero at 50\% --- one run each, and far short of uniform's 14. Uniform's
advantage comes from covering every span across 25 phases, not from any single proportional
remix.

\textbf{The same mask succeeds at the opposite task.} Following \citet{khodabandehlou2026abstention}, a LoRA trained on the same
free label to \emph{refuse} when evidence was evicted abstains on 0.728 of evicted probes against
0.000 for the base model (+17.5σ), cutting hallucination by 72.8\%, while refusing only 10.9\%
of probes whose answer is still present and raising retained accuracy from 0.684 to 0.770.
The compaction mask tells a model what was lost: enough to abstain, not enough to recover.

\textbf{Natural prose.} On HotpotQA (one seed, 160 probes, 92 evicted) nothing separates:
cascading with no adapter recovers 4 of 92, compaction 3, uniform 5 (+0.7σ uniform over
compaction). Some answers survive in other paragraphs or in the model's prior, so the
no-training floor is not zero. At this size the run neither replicates nor contradicts the
synthetic result.

\subsection{Gating distillation on the context gap}

\textbf{Scoring.} For a skill document d, demonstration query q and response r, form
\texttt{T\ =\ chat(system,\ d\ ++\ q)} and \texttt{S\ =\ chat(system,\ q)}, run the frozen backbone on \texttt{T\ ++\ r} and
\texttt{S\ ++\ r}, and score each response token by

\begin{verbatim}
kl_i = KL( p(· | T, r_<i) || p(· | S, r_<i) )
\end{verbatim}

Tokens group into sentence and code-line spans via the tokenizer's offset mapping; a span
scores the mean of its tokens.

\textbf{Gating and loss.} The top 25\% of spans get weight 1 and the rest a floor weight. The
teacher is the backbone with its adapter disabled and the document in context; the student
is the adapter-enabled model without it; the loss is \texttt{Σ\ w\_i\ KL\_i\ /\ Σ\ w\_i}. One model in
memory, one adapter per skill category.

\textbf{Setup.} Eight invented tool-use APIs in three categories, each with a \texttt{SKILL.md}, nine
demonstrations and twelve held-out tasks (96 in all); the APIs are fictional, so the base
model cannot know them. Qwen2.5-1.5B-Instruct, 300 steps at learning rate 1e-4. A task passes
when every required API string appears in the output. Arms: \texttt{uniform}, every token at full
weight (S2L-style \citep{zhang2026skill}); \texttt{random}, a random 25\% of spans at full weight and the floor
elsewhere; \texttt{kl\_top}, the KL-selected 25\% under the same floor.

\textbf{The gate carries a signal.} Required-token coverage --- the share of each demonstration's
required API identifiers inside the selection --- is \textbf{0.597 against 0.392} for a
matched-budget random selection, \textbf{+6.77σ}; at token granularity, 0.363 vs 0.247, +4.29σ.

\textbf{It does not help training.} At floor weight 0, \texttt{kl\_top} reaches 0.510 against \texttt{random}'s
0.583 and \texttt{uniform}'s 0.708: the prefix that conditions each selected identifier receives no
gradient. A floor of 0.1 repairs this:

\begin{longtable}[]{@{}llll@{}}
\toprule\noalign{}
seed & kl\_top & random & uniform \\
\midrule\noalign{}
\endhead
\bottomrule\noalign{}
\endlastfoot
0 & 0.729 & 0.729 & 0.708 \\
1 & 0.823 & 0.750 & 0.729 \\
2 & 0.802 & 0.719 & 0.698 \\
3 & 0.813 & 0.792 & 0.729 \\
4 & 0.656 & 0.729 & 0.677 \\
mean & 0.765 & 0.744 & 0.708 \\
\end{longtable}

\texttt{kl\_top} and \texttt{random} share identical masking, so their difference isolates the ranking:
\textbf{+0.021, t(4) = 0.74}, with the sign flipping between seeds. The ranking adds nothing over
random selection. Both masked arms beat full-weight \texttt{uniform} (\texttt{random} +0.035,
t(4) = 3.90; \texttt{kl\_top} +0.056, t(4) = 2.33, not significant). That is a separate finding about
masked versus full-weight training, and \texttt{uniform} is not simply mistuned: it gets
monotonically worse from 300 to 1,200 steps (0.708 to 0.656), and the best of five learning
rates from 5e-5 to 8e-4 reaches 0.750 (single seed each). At token granularity with floor 0.1,
\texttt{kl\_top} and \texttt{random} are again level (0.771 vs 0.760). The distilled skill keeps partial
function under a mismatched retrieved document (0.292 vs 0.000 for the prompted skill).

\subsection{Two failures, two reasons}

The compaction gate is \emph{harmful}: its selection systematically excludes whole fact types,
and no remixing tested restores them. The context-gap gate is \emph{neutral}: its selection is
informative about which tokens matter, but training on it is no better than training on a
random subset of the same size. Both signals clear matched controls; neither beats uncurated
coverage. A signal's strength says whether it carries information; it does not say whether
that information is what the weights are missing. Before gating training on any free signal,
compare it against uniform sampling at a matched budget --- the matched control of §6, one
level up.

\section{Matched controls, and what they caught}

A raw signal-strength figure is close to meaningless without a control beside it. We report
every one against a matched control because doing so repeatedly caught us:

\begin{longtable}[]{@{}
  >{\raggedright\arraybackslash}p{(\linewidth - 4\tabcolsep) * \real{0.3333}}
  >{\raggedright\arraybackslash}p{(\linewidth - 4\tabcolsep) * \real{0.3333}}
  >{\raggedright\arraybackslash}p{(\linewidth - 4\tabcolsep) * \real{0.3333}}@{}}
\toprule\noalign{}
\begin{minipage}[b]{\linewidth}\raggedright
Artifact
\end{minipage} & \begin{minipage}[b]{\linewidth}\raggedright
Apparent effect
\end{minipage} & \begin{minipage}[b]{\linewidth}\raggedright
Caught by
\end{minipage} \\
\midrule\noalign{}
\endhead
\bottomrule\noalign{}
\endlastfoot
Silent fallback to a rule-based compactor & 3.99x lift & decision-provenance counter \\
``Keep at most N'' answered \texttt{NONE} on every event & 100\% empty keeps & empty-keep rate \\
Model answered with the first N spans & 1.68x lift & positional control \\
\texttt{sorted(kept){[}:budget{]}} reimposed position after shuffling & 1.69x lift & kept-index distribution \\
Supporting paragraphs stacked at the front of a benchmark & 1.76x lift & positional control (2.44x) \\
Cached scores from an older tokenizer inverted the gate & −5.14σ & re-scoring in the current environment (+6.40σ) \\
Falling CE on gold answers read as knowledge transfer & 6.29 → 2.24 & distractor ranking (at chance) \\
Eval questions that did not identify which item was meant & 50\% ceiling & duplicate-answer audit \\
\end{longtable}

Six of the eight inflated a result in our favour; two hid a possible positive. None was a
modelling error; all were construction, which we think is the norm.

\textbf{The tokenizer row} moved the result against us, so disbelief did not catch it, and we
misdiagnosed it once. Scored from a cached file produced under an older tokenizer, which
split identifiers such as \texttt{vx\_stash} into five pieces, code spans averaged low and the gate
scored 0.118 required-token coverage against a matched control's 0.307. We blamed
mean-pooling. Re-scoring in the current environment, with no code change, gave 0.586 against
0.392. Having seen the stale number, we also marked the downstream comparison confounded;
re-running it reproduced the original result exactly. A retraction needs the same evidence
as a claim, and a cached score file is valid only in the environment that produced it.

\textbf{The last two rows compound.} Consolidation lowered per-token CE on evicted gold answers
from 6.29 to 2.24, which we first read as knowledge that had entered the weights; ranking
each gold answer against same-bank distractors put every arm at the 0.240 chance rate, which
we then read as knowledge never acquired. Both were premature. 98.3\% of the benchmark's
questions had more than one correct answer across trajectories --- ``Which header carries the
auth token?'' was asked of 48 conversations with four different answers --- so a perfect
memoriser could not exceed 50\%, and the distractors were other trajectories' correct answers.
Scoping every question to its trajectory is what made §5.1 measurable. The audit --- count the
distinct answers each question string receives --- costs minutes and should precede the
experiment.

\section{Discussion}

The practical rule for free-supervision pipelines below frontier scale is to \textbf{derive the
signal from what the model does, not from what it says it does}, and then to check the
signal against uniform sampling before training on it. The context gap obeys the first half
by construction; the compaction decision obeys it once read rather than generated;
Declarative Attention as published does not, and §4.2 measures the cost at small scale. The
second half is where both consolidation signals fail.

The two halves interact. The measured compaction decision is strong precisely because it
reflects what the compactor considers salient --- and that is also why it fails as a training
target, since what a compactor keeps for immediate use need not be what the weights lack. A
signal that is maximally informative about one decision can be systematically biased for
another. This suggests the useful free signals for consolidation are ones whose errors are
uncorrelated with coverage, or ones used, as in abstention, for the decision they were made
for.

For routing, the logit read is a cheap, training-free substitute for declaring at the scales
where declaring is unreliable. It is not a deployable router on its own --- about 3x chance
leaves most probes misrouted --- and it costs K region prompts rather than a single generated
token. Whether it keeps its lead where declaration works, at 27B and above, is the natural
next test.

\section{Limitations}

\begin{itemize}
\tightlist
\item
  \textbf{Synthetic consolidation.} The significant consolidation result is on a controlled
  synthetic benchmark. The HotpotQA run is one seed at 92 evicted probes with a non-zero
  no-training floor and cannot resolve a gap of the synthetic size. The synthetic generator's
  unmarked variant is a floor case: facts and filler share one template bank.
\item
  \textbf{Seeds.} The main comparisons have five seeds (§5.1 corpora, §5.2 gate) or three (§4.2);
  the remix ablations, abstention, the step and learning-rate sweeps, and HotpotQA are single
  runs.
\item
  \textbf{Scale.} Everything is at or below 7B, and 7B is 4-bit quantised. We cannot test the
  27B+ regime of \citet{ho2026attention}, so our result bounds the generated declaration from below --- at a
  fixed-slot control at 1.5B, under half the logit read's signal at 7B --- and does not
  contradict the accuracy \citet{ho2026attention} reports at scale. The logit read is untested there.
\item
  \textbf{Routing accuracy.} The best router is about 3x chance; §4.2 compares elicitations, it
  does not deliver a usable router.
\item
  \textbf{One model family.} Apart from SmolLM2-360M in §4.1, all results are on Qwen2.5.
\end{itemize}

\section{Conclusion}

Free salience signals exist at small scale, but only when measured: asked to state what
matters, 0.5B--7B models fail or carry a fraction of the signal their logits hold. And a
measured signal that clears its control can still be the wrong thing to train on ---
compaction-gated consolidation loses to uniform replay on every corpus, and context-gap
gating adds nothing over random selection. Both lessons reduce to the same practice: put a
matched control beside every number, including the number that says a signal should be
trained on.

\bibliography{main}
\bibliographystyle{tmlr}

\newpage
\appendix

\section{Hyperparameters}

\begin{longtable}[]{@{}
  >{\raggedright\arraybackslash}p{(\linewidth - 6\tabcolsep) * \real{0.2500}}
  >{\raggedright\arraybackslash}p{(\linewidth - 6\tabcolsep) * \real{0.2500}}
  >{\raggedright\arraybackslash}p{(\linewidth - 6\tabcolsep) * \real{0.2500}}
  >{\raggedright\arraybackslash}p{(\linewidth - 6\tabcolsep) * \real{0.2500}}@{}}
\toprule\noalign{}
\begin{minipage}[b]{\linewidth}\raggedright
\end{minipage} & \begin{minipage}[b]{\linewidth}\raggedright
§5.1 consolidation
\end{minipage} & \begin{minipage}[b]{\linewidth}\raggedright
§5.2 distillation
\end{minipage} & \begin{minipage}[b]{\linewidth}\raggedright
§4.2 routing
\end{minipage} \\
\midrule\noalign{}
\endhead
\bottomrule\noalign{}
\endlastfoot
backbone & Qwen2.5-0.5B (compactor 1.5B) & Qwen2.5-1.5B & Qwen2.5-0.5B / 1.5B / 7B \\
LoRA & r 16, α 32, dropout 0.05, all projections & same & --- \\
learning rate & 5e-5 & 1e-4 & --- \\
steps & 80 per sleep phase & 300 & --- \\
schedule & sleep every 4 compaction events, 25 phases & one adapter per category & --- \\
replay & reservoir, capacity 80 & --- & --- \\
selection & keep ρ = 0.25, sentence spans, 8 recent turns protected & top 25\% of spans, floor 0 or 0.1 & K = 8 regions \\
data & 48 trajectories × 120 turns × 6 facts & 8 APIs, 9 demos, 12 tasks each & 96 HotpotQA trajectories × 4 probes \\
seeds & 5 corpora & 5 & 3 layouts \\
precision & fp16 & fp16 & fp16; 7B nf4 \\
\end{longtable}

\section{Reproducibility}

All code, configs and Kaggle notebooks are in the supplementary repository. Each result
above maps to one config and one notebook; \texttt{docs/project\_a\_findings.md},
\texttt{docs/project\_b\_findings.md} and \texttt{docs/declarative.md} hold the per-run tables, and
\texttt{eval/declare\_report.py} regenerates the §4.2 table and paired tests from the run
directories. The 7B routing runs use \texttt{configs/kaggle\_declare\_7b.yaml} with a repeat-KV SDPA
attention (\texttt{declare/elicit.py}) that matches stock SDPA exactly but avoids a 5 GB attention
matrix on T4 hardware for 9k-token prompts.

\end{document}
